\section{Steady-State Response and Numerical validation}
\subsection{Simulation setup and observables}
\quad We perform Langevin simulations on synthetic networks to validate the theoretical response functions. 
Both undirected and directed networks are considered, with either uniform or heterogeneous noise. The network dynamics follow (\ref{noisy_network}), and the steady-state covariance matrix $\mathbf{K}_0 = \langle \delta \vec{x} \delta \vec{x}^\top \rangle$ is measured from long-time trajectories with time step $\Delta t = 0.0005$ and simulation period $T=10000$. To probe the system response, we introduce small perturbations to either node parameters $\delta r_\alpha$ or connection weights $\delta W_{\alpha\beta}$, where the perturbations of multiple nodes and edges are also considered to test the robustness, and these can be performed by small modifications to parameters. Disruptions of networks can be realized bu decreasing the connection weights of chosen edges to zero. For disruptions of nodes, they become isolated nodes where all edges incident on them are removed (connection weights of all edges connected with the chosen nodes become zero). After perturbation, the system relaxes to a new steady state, from which the updated covariance matrix is measured. The response is quantified by the difference between the perturbed and unperturbed steady states, $\Delta \mathbf{K}_0$, and compared with the theoretical prediction obtained from the unperturbed steady-state time series. We consider the undirected/directed ER random graphs with number of nodes $N=50, 100$ and the connection probability is about $0.1\sim0.2$. The intrinsic dynamics of each node is governed by the Logistic function $f_i=r_ix_i(1-x_i)$ and the interaction between nodes is the synchronizing coupling function $h(x_i,x_j)=x_j-x_i$. Under this condition, the Jacobian matrix $\mathbf{Q}$ is given by
\begin{equation}
    Q_{ii} = -r_i - \sum_{j\neq i} W_{ij},\; Q_{ij} = W_{ij},
\end{equation}
where the noise-free fixed points for each node is $1$. The network parameters are chosen to represent generic noisy network dynamics while allowing controlled comparisons between equilibrium and nonequilibrium steady states. From the detailed balance condition $\mathbf{Q}\boldsymbol{\sigma}=\boldsymbol{\sigma}\mathbf{Q}^\intercal$, an asymmetric $\mathbf{W}$ with a uniform diagonal $\boldsymbol{\sigma} = \sigma_0 \mathbf{I}$ and a symmtric $\mathbf{W}$ with a nonuniform diagonal $\boldsymbol{\sigma}$ always break detailed balance. The equilibrium case is the undirected graph with uniform noise strengths or the special choise $W_{ij} \sigma_{jj} = W_{ji} \sigma_{ii}$. As for the intrinsic node parameter, which does not affect detailed balance, are only chosen to ensure the existence of a stable fixed point. We consider the non-negative $r_i$ and $W_{ij}$ and the noise strenghs of all nodes are positive.\par
While the simulations on smaller networks clearly demonstrate the validity of the response formulas, 
it remains important to verify that the agreement is not restricted to finite-size or realization-specific effects. From a computational perspective, large networks are particularly relevant because the covariance matrix can be obtained efficiently from the Lyapunov equation, allowing the response functions to be evaluated without explicitly repeating perturbation simulations for every parameter change.Since the present framework is intended as a general response theory for noisy network dynamcis around fixed points, we also consider networks with a larger number of nodes and different classes of network topologies, including regular lattices, WS small-world networks, and scale-free networks. These network models represent distinct structural features, ranging from homogeneous local connectivity to highly heterogeneous degree distributions.

\subsection{Equilibrium response: undirected networks with uniform noise}
\quad We first consider undirected networks with uniform noise strengths, for which the steady state satisfies detailed balance. We perturb either node parameters $\delta r_\alpha$ or connection weights $\delta W_{\alpha\beta}$, and measure the resulting change in the steady-state covariance matrix $\mathbf{K}_0$. The theoretical prediction is obtained from the unperturbed steady-state statistics. Fig. \ref{chieqmlinear}a and b show the comparison between the simulation of the change in $\mathbf{K}_0$ and the theoretical response function with respect to the perturbation of a single node and edge, the response function is given as (\ref{chiKodr_eq}) and (\ref{chiKodW_eq}). The node parameters are uniformly distributed on $[0,5]$, and the connection weights are normally distributed $W_{ij} \sim \mathcal{N}(5,1)$. The noise strengths are identical among the nodes $\sigma_{ii}=0.01$. The data points collapse on the diagonal, which indicates the validation of the response function for noisy network dynamics at equilibrium for both node and edge perturbations. In Fig. \ref{chieqmlinear}a, the perturbed node ($\alpha$) is randomly chosen with three different positive perturbation values $dr_\alpha = 0.1, 0.5, 2$, denoted as circles, squares, and triangles. The node parameters determine the strength of the intrinsic forces on the nodes, so the increase of r stabilizes the nodes and reduces the effect of noise, and the elements in $\mathbf{K}_0$ decrease. The change in $(\mathbf{K}_0)_{ij}$ depends on the distance between the perturbed node $\alpha$ and the other nodes. The data point at the bottom left corner is $(\mathbf{K}_0)_{\alpha\alpha}$, which shows the perturbed node has the largest response. The change in $(\mathbf{K}_0)_{i\alpha}$ where node $i$ and the perturbed one has a direct connection $W_{i\alpha}$ and $W_{\alpha i} > 0$ is around $-5e-7$, and the other nodes that do not connect to $\alpha$ hardly respond to the perturbation.
\begin{figure}[H]
    \centering
    \subfigure[]{\includegraphics*[width=.45\columnwidth]{Figure/Steady/Undirected/dKdr.eps}}
    \subfigure[]{\includegraphics*[width=.45\columnwidth]{Figure/Steady/Undirected/dKdW.eps}}
    \caption{The two figures show the Langevin dynamics simulations for an undirected ER random network of $N=50$ nodes with linear force intrinsic node dynamics $f_i(x_i)=-r_i(x_i-1)$ and connection probability $p=0.2$ with a uniform noise matrix a normally distributed $\bf{W}$: (a) The figure shows the response function of $\bf{K}_0$ due to the change in $r_\alpha$ and $d r_\alpha = 0.1, 0.5, 2$ and is plotted against the simulation results. $\{r_i\}$ is normally distributed with the mean value $\overline{r}=5$. The point at the bottom left corner is the response of $({\bf K}_0)_{\alpha\alpha}$, which has the largest response. The points around $-5e-7$ are the nodes directly connected to $\alpha$, and those around 0 have no path to node $\alpha$. (b) The figure shows the response function of ${\bf K}_0$ due to the change in $W_{\alpha\beta}$, where the mean value $\overline{W} = 5$. The three points from bottom left to top right are the response of $({\bf K}_0)_{\beta\beta}$, $({\bf K}_0)_{\alpha\alpha}$, and $({\bf K}_0)_{\alpha\beta}$ respectively.}
    \label{chieqmlinear} 
\end{figure}
In Fig. \ref{chieqmlinear}b, the perturbed edge $(\alpha,\beta)$ is randomly chosen with three different positive perturbation values $dW_{\alpha\beta} = 0.2, 1, 5$, denoted as circles, squares, and triangles. Notice that for undrected networks, the $dW_{\beta\alpha}$ also changes to maintain the symmetric topology. The coupling strengths between any pair of nodes are characterized by the connection weights. The increase of the connection weight between node $\alpha$ and node $\beta$ enhances the correlation between the states of the two nodes. The data point at the top right corner is $(\mathbf{K}_0)_{\alpha\beta}$, and the auto-correlations $(\mathbf{K}_0)_{\alpha\alpha}$ and $(\mathbf{K}_0)_{\beta\beta}$ are at the left bottom corner, where the increasing coupling reduces the fluctuations of the two nodes. In the case of the Logistic force function, the response function of $\bf{K}_0$ is the same as that in the linear force function. As Fig. \ref{chieqmlogistic}a and \ref{chieqmlogistic}c show, the theoretical predictions match the simulation results when a parameter $r_\alpha$ or $W_{\alpha\beta}$ changes.
\begin{figure}[H]
    \centering
    \subfigure[]{\includegraphics*[width=.45\columnwidth]{Figure/Steady/Undirected/chiij1eqm.eps}}
    \subfigure[]{\includegraphics*[width=.45\columnwidth]{Figure/Steady/Undirected/LogisticKowun.eps}}
    \caption{Langevin dynamics simulations for an undirected ER random network of $N=50$ nodes with heterogeneous Logistic intrinsic node dynamics $f_i(x_i)=r_ix_i(1-x_i)$ and connection probability $p=0.2$, with a uniform diagonal noise matrix. (a) The response function of the fluctuation correlations. The change in the elements of the correlation functions upon a change of the node parameter of $\delta r_\alpha$ of the node $\alpha$ is measured by simulations. Four nodes are arbitrarily chosen with three different values of $\delta r_\alpha$. $\chi_{ij\alpha}$ is significantly larger for $i=j=\alpha$. (b) The response of ${\bf K}_0$ due to the change in $\bf W$ is large at $(i,j)=(\alpha,\alpha)$, $(\alpha,\beta)$, and $(\beta,\beta)$.}
    \label{chieqmlogistic} 
\end{figure}

\subsection{Nonequilibrium steady-state response}
\quad  In our framework, a NESS arises when detailed balance is broken, which can occur either through asymmetric interactions (directed networks) or heterogeneous noise distributions. In contrast to equilibrium systems, where the steady state is characterized by vanishing probability currents, NESS exhibit persistent irreversible fluxes. The key question we address is whether the steady-state linear response framework remains predictive under such conditions. To highlight the difference between equilibrium and nonequilibrium responses, we compare the exact response with the equilibrium-like simplified formula of $\mathbf{K}_0$, as given in (\ref{chiKodr_eq}) and (\ref{chiKodW_eq}). In Fig. \ref{chieqcompare}, a single node and a single directed edge are perturbed by $\delta r_\alpha, \delta W_{\alpha\beta}=0.5$ in a directed ER random network with a uniform noise matrix. The vertical axis is the equilibrium response function $\chi^{eq}_{ij\alpha(\beta)}$ and the horizontal one is the exact response function $\chi^{ss}_{ij\alpha(\beta)}$ solved from the Lyapunov equation. 
\begin{figure}[H]
    \centering
    \subfigure[]{\includegraphics*[width=.4\columnwidth]{Figure/Else/chieqnondr.eps}}
    \subfigure[]{\includegraphics*[width=.4\columnwidth]{Figure/Else/chieqnondW.eps}}
    \subfigure[]{\includegraphics*[width=.4\columnwidth]{Figure/Else/chieqdirdr.eps}}
    \subfigure[]{\includegraphics*[width=.4\columnwidth]{Figure/Else/chieqdirdW.eps}}    
    \caption{The comparison between the equilibrium and the NESS response functions for the covariance matrix due to the perturbations of a single node and an edge: When the detailed balance condition is violated, the response must instead be computed from the full Lyapunov response equation.}
    \label{chieqcompare} 
\end{figure}
Similar to the equilibrium case, we compare the simulation results with the theoretical response functions for the NESS conditions: a directed graph with a uniform diagonal noise matrix in Fig. \ref{chiKodir} and an undirected graph with a nonuniform diagonal noise matrix in Fig. \ref{chiKonon}. In both two cases, the perturbed node and the perturbed edge are randomly chosen, and the perturbed values are $\delta r_\alpha = 0.5, 2$ and $\delta W_{\alpha\beta} = 0.5, 2$. The distributions of the network parameters are the same as the equilibrium cases in Fig. \ref{chieqmlinear} and Fig. \ref{chieqmlogistic}. As the figures show, the response functions of the equal-time covariance are valid for the directed graph and the heterogenous noise strengths, where detailed balance is broken. Note that in the directed graph, we apply a single edge perturbation; however, in the undirected graph, when one edge $W_{\alpha\beta}$ is perturbed, the other edge $W_{\beta\alpha}$ with the opposite direction is also perturbed with the same amount to maintain the undirected topology after the perturbation. In Fig. \ref{chiKodir}b, the perturbed edge is from node $\beta$ to node $\alpha$, and thus only node $\alpha$ is affected. There is a data point $\delta (\mathbf{K}_0)_{\alpha\alpha} / \delta W_{\alpha\beta}$ at the bottom left corner. In Fig. \ref{chiKonon}b, both the reciprocal edges are perturbed with the same amount of the perturbation value, so node $\alpha$ and $\beta$ are perturbed. There are two data points at the bottom left corner, which indicates the response of $(\mathbf{K}_0)_{\alpha\alpha}$ and $(\mathbf{K}_0)_{\beta\beta}$ decrease due to the increase of the coupling between them.
\begin{figure}[H]
    \centering
    \subfigure[]{\includegraphics*[width=.45\columnwidth]{Figure/Steady/Directed/dKdr_dir.eps}}
    \subfigure[]{\includegraphics*[width=.43\columnwidth]{Figure/Steady/Directed/dKdW_dir.eps}}
    \caption{Comparison between Langevin simulations and the steady-state response prediction for a directed ER random network with a uniform diagonal noise matrix, where detailed balance is broken by asymmetric couplings. (a) Response of $\mathbf{K}_0$ to perturbations of a single node parameter with $\delta r_\alpha = 0.5, 2$. (b) Response of $\mathbf{K}_0$ to perturbations of a single directed edge weight with $\delta W_{\alpha\beta} = 0.5, 2$. The collapse of the data points onto the diagonal confirms the validity of the NESS response function obtained from the Lyapunov equation.}
    \label{chiKodir} 
\end{figure}

\begin{figure}[H]
    \centering
    \subfigure[]{\includegraphics*[width=.45\columnwidth]{Figure/Steady/Undirected/dKdr_non.eps}}
    \subfigure[]{\includegraphics*[width=.45\columnwidth]{Figure/Steady/Undirected/dKdW_non.eps}}
    \caption{Comparison between Langevin simulations and the steady-state response prediction for an undirected ER random network with a nonuniform diagonal noise matrix, where detailed balance is broken by heterogeneous noise strengths. (a) Response of $\mathbf{K}_0$ to perturbations of a single node parameter with $\delta r_\alpha = 0.5, 2$. (b) Response of $\mathbf{K}_0$ to perturbations of a single undirected edge, where both $W_{\alpha\beta}$ and $W_{\beta\alpha}$ are increased by the same amount $\delta W_{\alpha\beta} = 0.5, 2$ to preserve the symmetric topology. The agreement with the diagonal validates the NESS response function in the presence of heterogeneous noise.}
    \label{chiKonon} 
\end{figure}

\subsection{Multiple perturbations and removal of network components}
\quad After validating the linear-response theory for single-parameter perturbations, we next examine whether the response functions can be superposed to predict the effect of multiple simultaneous perturbations. We define the covariance change as
\begin{equation}
    \Delta (\mathbf{K}_0)_{ij} =
    (\mathbf{K}_0^{\rm after})_{ij}-(\mathbf{K}_0^{\rm before})_{ij}.
\end{equation}
For a set of perturbed nodes
\(\mathcal{P}_V=\{\alpha_1,\alpha_2,\ldots,\alpha_m\}\), the first-order prediction is
\begin{equation}
    \Delta (\mathbf{K}_0)_{ij}
    =
    \sum_{\alpha\in\mathcal{P}_V}
    \chi^{(r)}_{ij\alpha}\delta r_\alpha .
\end{equation}
Similarly, for a set of perturbed edges \(\mathcal{P}_E=\{(\alpha_1,\beta_1),(\alpha_2,\beta_2),\ldots, (\alpha_m,\beta_m)\}\), we use
\begin{equation}
    \Delta (\mathbf{K}_0)_{ij}
    =
    \sum_{(\alpha,\beta)\in\mathcal{P}_E}
    \chi^{(W)}_{ij(\alpha\beta)}\delta W_{\alpha\beta}.
\end{equation}
The perturbation amplitudes \(\delta r_\alpha\) and \(\delta W_{\alpha\beta}\) are allowed to be different for different perturbed nodes or edges. To test this superposition property, we perform Langevin simulations on undirected and directed ER networks with nonuniform noise sources.  In each realization, multiple nodes or edges are randomly selected and perturbed simultaneously. The covariance change measured from the simulations is then compared with the first-order prediction obtained by summing the corresponding single-parameter response functions.  Agreement with the diagonal line in the scatter plots indicates that the covariance response remains additive within the linear regime.\par
In Fig. \ref{MultipledKdr}, ten randomly selected nodes are perturbed with the same amount of perturbation values $\delta r_\alpha=0.5$, where the set of the perturbed nodes is denoted as $\mathcal{P}_V=\{ \alpha_1, \alpha_2, ..., \alpha_{10} \}$. For visual clarity, we separate the elements of the covariance matrix into diagonal and off-diagonal components. We also classify the elements according to whether their indices include nodes associated with the perturbation: diagonal elements are labeled as $\alpha\alpha$ or $ii,\;(i\notin \mathcal{P}_V)$, and off-diagonal elements as $i\alpha$ or $ij,\;(i,j \notin \mathcal{P}_V)$. The data points collapse onto the diagonal line, indicating quantitative agreement between the linear-response prediction and the numerical simulation for multiple node perturbations. In Fig. \ref{MultipledKdW}, twenty randomly selected undirected edges are perturbed with the same amount of values $\delta W_{\alpha\beta}=0.5$. The perturbed edges are represented as the set $\mathcal{P}_E = \{ (\alpha_k,\beta_k),\;k=1,2,...,20 \}$. Again, the data points are separated into the diagonal and offdiagonal elements. The elements are further classified as $\alpha\alpha$ or $ii,\;(i\neq \alpha)$ for the diagonal ones, and $i\alpha$ or $ij,\;(i,j\neq \alpha)$. Notice that the classification of edge perturbations is based on the receiving node $\alpha$, not on the source node $\beta$. The data points align along the diagonal, which proves the validation of the linear response framework on multiple edge perturbations.
\begin{figure}[H]
    \centering
    \subfigure[]{\includegraphics*[width=.45\columnwidth]{Figure/Multiple/MultipledKdrii_non.eps}}
    \subfigure[]{\includegraphics*[width=.45\columnwidth]{Figure/Multiple/MultipledKdrij_non.eps}}
    \caption{Comparison between Langevin simulations and the superposed linear-response prediction for simultaneous perturbations of ten node parameters in an ER random network with nonuniform noise. (a) Diagonal covariance changes $\Delta(\mathbf{K}_0)_{ii}$. (b) Off-diagonal covariance changes $\Delta(\mathbf{K}_0)_{ij}$. The grouping of data points distinguishes entries involving perturbed nodes from those that do not, and the collapse onto the diagonal confirms the additivity of the first-order response.}
    \label{MultipledKdr} 
\end{figure}
\begin{figure}[H]
    \centering
    \subfigure[]{\includegraphics*[width=.44\columnwidth]{Figure/Multiple/MultipledKdWii_non.eps}}
    \subfigure[]{\includegraphics*[width=.46\columnwidth]{Figure/Multiple/MultipledKdWij_non.eps}}
    \caption{Comparison between Langevin simulations and the superposed linear-response prediction for simultaneous perturbations of twenty undirected edges in an ER random network with nonuniform noise. (a) Diagonal covariance changes $\Delta(\mathbf{K}_0)_{ii}$. (b) Off-diagonal covariance changes $\Delta(\mathbf{K}_0)_{ij}$. The alignment of the data along the diagonal shows that the first-order responses to individual edge perturbations remain additive in the linear regime.}
    \label{MultipledKdW} 
\end{figure}
\begin{figure}[H]
    \centering
    \subfigure[]{\includegraphics*[width=.42\columnwidth]{Figure/Multiple/MultipledKdWii_remove.eps}}
    \subfigure[]{\includegraphics*[width=.46\columnwidth]{Figure/Multiple/MultipledKdWij_remove.eps}}
    \caption{Comparison between Langevin simulations and the first-order prediction for the removal of 40 directed edges in a directed ER random network with nonuniform noise. The theoretical prediction is obtained by summing the single-edge response functions with $\delta W_{\alpha\beta}=-W_{\alpha\beta}$ for all removed edges. (a) Diagonal covariance changes $\Delta(\mathbf{K}_0)_{ii}$. (b) Off-diagonal covariance changes $\Delta(\mathbf{K}_0)_{ij}$. The approximate collapse onto the diagonal shows that the linear response still captures the dominant effect of multiple edge removals.}
    \label{MultipledKdW_remove}
\end{figure}
We further consider removal events as a finite-amplitude stress test of the linear-response framework. Edge removal is implemented by setting \(W_{\alpha\beta}\to 0\), corresponding to \(\delta W_{\alpha\beta}=-W_{\alpha\beta}\). Node removal is implemented by isolating the selected node through the removal of all incoming and outgoing edges connected to it. Since these perturbations are not infinitesimal, the first-order prediction is not expected to be quantitatively exact. Instead, this comparison identifies the regime in which the linear response remains informative and the regime in which higher-order effects become important. Fig. \ref{MultipledKdW_remove} examines the removal of 40 directed edges in a directed ER network with nonuniform noise. The prediction is obtained by summing the single-edge response functions over all removed edges. The agreement with the diagonal line shows that, for this level of edge removal, the first-order response still captures the dominant covariance change, even in a nonequilibrium steady state generated by directed couplings and heterogeneous noise. Figure \ref{MultipledKdW_noderemove} considers the removal of a single node in an undirected network with nonuniform noise. Although only one node is removed, this operation simultaneously deletes all edges connected to that node and therefore represents a large structural perturbation. The visible deviations from the diagonal line are expected, especially for covariance elements involving the removed node or its neighbors. These deviations indicate the onset of higher-order response terms beyond the first-order response theory.
\begin{figure}[H]
    \centering
    \subfigure[]{\includegraphics*[width=.45\columnwidth]{Figure/Multiple/MultipledKdWii_noderemove.eps}}
    \subfigure[]{\includegraphics*[width=.45\columnwidth]{Figure/Multiple/MultipledKdWij_noderemove.eps}}
    \caption{Comparison between Langevin simulations and the first-order prediction for the removal of a single node in an undirected ER random network with nonuniform noise. The prediction is obtained by summing the responses of all incident edge removals associated with the deleted node. (a) Diagonal covariance changes $\Delta(\mathbf{K}_0)_{ii}$. (b) Off-diagonal covariance changes $\Delta(\mathbf{K}_0)_{ij}$. The visible deviations from the diagonal indicate the onset of higher-order effects for this large structural perturbation.}
    \label{MultipledKdW_noderemove}
\end{figure}
% \subsection{Large-scale networks and topology dependence}
% \quad To assess the scalability and topological applicability of the linear response framework, we next consider large-scale networks with $N > 1000$. In this regime, direct Langevin simulations become computationally expensive for estimating steady-state $\mathbf{K}_0$. We therefore validate the linear response predictions by comparing them with covariance changes obtained from direct solutions of the Lyapunov equation before and after parameter perturbations. This procedure tests whether the first-order response formula remains quantitatively accurate in large networks and whether its performance depends on network topology. We consider the synthetic networks such as Erdős-Rényi, Barabási-Albert, Watts-Strogatz, and ring lattices, which represent random, scale-free, small-world, and local nearest-neighbor connectivity, respectively. In Fig. \ref{Response_for_large_undirected_network}, multiple perturbations of nodes and edges are applied in the four types of undirected network topologies with $N=1000$. 40 randomly selected nodes are perturbed with the perturbation values on the intercal ($\delta r$ in $[-1, 1]$), and 100 undirected edges are perturbed with the perturbation values on the intercal ($\delta W$ in $[-1, 1]$). The mean values of the node parameters $r_i$ and connection weights $W_{ij}$ are fixed at $5$ and they are sampled from a uniform distribution on $[2.5, 7.5]$ and a Gaussian distribution $\mathcal{N}(5,0.5)$. The data points collapse on the diagonal for the comparison between the analytic solution of $\mathbf{K}_0$ from the Lyapunov equation and the response functions, these results show that the linear response framework and the superposition are also valid for the linearized noisy network dynamics with different topologies.

% \begin{figure}[H]
%     \centering
%     \subfigure[]{\includegraphics*[width=.45\columnwidth]{Figure/Large_network/ER_BA_undir_dr.eps}}
%     \subfigure[]{\includegraphics*[width=.45\columnwidth]{Figure/Large_network/WS_RING_undir_dr.eps}}
%     \subfigure[]{\includegraphics*[width=.47\columnwidth]{Figure/Large_network/ER_BA_undir_dW.eps}}
%     \subfigure[]{\includegraphics*[width=.45\columnwidth]{Figure/Large_network/WS_RING_undir_dW.eps}}
%     \caption{The }
%     \label{Response_for_large_undirected_network}
% \end{figure}

% After verifying the quantitative agreement between the theoretical response functions and numerical measurements, we next ask how the magnitude of the response is controlled by network structure. In particular, we focus on the local variance response $\Delta (\mathbf{K}_0)_{\alpha\alpha}$ and the pairwise covariance response $\Delta (\mathbf{K}_0)_{\alpha\beta}$. The former quantifies how the fluctuation amplitude of the perturbed or directly affected node changes, and therefore provides a measure of local sensitivity or robustness. The latter quantifies how the dynamical correlation between two nodes is modified by a node or edge perturbation, and therefore captures how local perturbations are transmitted through network interactions. Together, these two quantities separate the self-response of a node from the correlation response between coupled nodes.
